\chapter{Yang--Mills 梯度流纲领}

2009--2010 年间，Martin L\"uscher 为规范场引入一个扩散方程，并将其
锻造成现代格点场论中最富成果的思想之一。流在半径 $\sqrt{8t}$ 上平滑
规范场且保持规范不变性，而光滑化后的场是\emph{重整化的场}：由其构成
的复合观测量无需附加重整化即有限。本库收藏了该纲领的原始论文；它们
支撑着两大分支——一侧是精密标度设定与热力学，另一侧是涂抹算符乘积
展开与涂抹准 PDF。

% ----------------------------------------------------------------------
\section{光滑 Wilson 圈：流之前的流思想}
\papercard{R.\ Narayanan and H.\ Neuberger}
{连续 Wilson 圈算符中的无穷 N 相变}
{JHEP \textbf{03} (2006) 064}
{arXiv:hep-th/0601210 | DOI: 10.1088/1126-6708/2006/03/064}
\whyselected{本库原文之一，也是 Wilson 线梯度流涂抹的概念祖先：
Monahan--Orginos 与 Monahan 的准 PDF 论文、Brambilla--Wang 的离光锥
论文以及高阶流论文都引用它（约 $\sim$6/99）。}
\physics{大尺寸 Wilson 圈在连续极限下并非良定义观测量——其周长发散。
Narayanan 与 Neuberger 研究了通过扩散轮廓得到的\emph{光滑化}圈，
并数值证明光滑圈具有有限的连续极限。在 't~Hooft 极限下，放大圈会触发
一个相变：和乐群的本征值密度闭合谱隙；若该相变属于可解普适类
（文中引用的 Gross--Witten / Durhuus--Olesen / Douglas--Kazakov 谱系），
弦张力便可解析触及。}
\impact{超出大 $N$ 物理本身，这篇论文确立了一个事实：对 Wilson 线做
光滑化与有限连续极限相容——这正是涂抹（伪/准）PDF 的立足点，因为其
非局域算符含开 Wilson 线。库内的流式 PDF 论文均以它为思想源头引用。}
\cardend

% ----------------------------------------------------------------------
\section{平庸化映射：流进入场论}
\papercard{M.\ L\"uscher}{平庸化映射、Wilson 流与 HMC 算法}
{Commun.\ Math.\ Phys.\ \textbf{293}, 899--911 (2010)}
{arXiv:0907.5491 | DOI: 10.1007/s00220-009-0995-5}
\whyselected{本库原文之一，文献中 Yang--Mills 梯度流的首次登场：
"Wilson 流"在此得名。库内流语料普遍引用（约 $\sim$6/99）。}
\physics{平庸化映射把作用量变换为长程耦合消失的作用量，理想情况下
使测度高斯化。L\"uscher 证明其无穷小版本导致场空间中的扩散方程
\begin{equation}\label{eq:trivmap}
  \partial_s B_\mu(s,u) = E_\mu(B(s,u)), \qquad B_\mu(0)=U_\mu ,
\end{equation}
其中 $E_\mu$ 是作用量在 $SU(3)$ 上的负梯度——即流。精确的映射会使
蒙特卡洛变得平庸；近似的映射仍可用于 HMC 内部以抑制拓扑自关联。
文章已指出短流时间探测理论的紫外结构，为后续一切工作设定了议程。}
\impact{正是这篇文章让共同体认识到：流的场既是数值工具又是重整化
概念。第八章的 AI 采样论文在严格意义上就是对可学习的平庸化映射之
搜索。}
\cardend

% ----------------------------------------------------------------------
\section{Wilson 流的性质}
\papercard{M.\ L\"uscher}{格点 QCD 中 Wilson 流的性质与应用}
{JHEP \textbf{08} (2010) 071}
{arXiv:1006.4518 | DOI: 10.1007/JHEP08(2010)071}
\whyselected{本次挖掘的全部引文网络中被引最高的单篇（约 $\sim$15/99）：
库内每篇梯度流论文都建立在其上，连 LaMET 与 BMW 论文也为其标度设定
或 QED 流插值而引用。本库原文之一。}
\physics{Wilson 流让链场按规范协变热方程演化
\begin{equation}\label{eq:flow}
  \partial_t B_\mu = D_\nu G_{\nu\mu} + \xi\,D_\mu\,\partial_\nu B_\nu,
  \qquad B_\mu(0)=A_\mu ,
\end{equation}
$\xi=0$（Wilson 方案）或 $1/5$（对称方案）。结合单圈微扰与数值证据，
L\"uscher 确立三个事实：(i)~任意固定 $t>0$ 下 $B_\mu(t)$ 是光滑场，
其局域规范不变泛函具有有限连续极限——流本身就是一种重整化；
(ii)~观测量 $w(t)=t^2\langle E(t)\rangle$
（$E=\frac12 G^a_{\mu\nu}G^a_{\mu\nu}$）在 $t\sim a^2$ 处陡峭穿越
crossover，故由 $\sqrt{8t}\,w(t)|_{w_0}=w_0$
定义的标度 $w_0$（约 $0.4$ fm）提供精密参考标尺；
(iii)~一旦流动超过 $\sqrt{8t}\sim a$，拓扑扇区动力学退耦。}
\impact{$w_0/w_a$ 标度设定已成为库内各系综（BMW、CLQCD clover 组）
的标准做法。更深刻的是"流的乘积紫外有限"这一定理——它是本章及第六章
所有涂抹算符应用的许可证。}
\cardend

% ----------------------------------------------------------------------
\section{流理论的可重整性}
\papercard{M.\ L\"uscher and P.\ Weisz}{非阿贝尔规范理论中梯度流的微扰分析}
{JHEP \textbf{02} (2011) 051}
{arXiv:1101.0963 | DOI: 10.1007/JHEP02(2011)051}
\whyselected{流语料第二大枢纽（约 $\sim$13/99）：库内全部短流时间计算
（NNLO 能动张量、NNLO 夸克双线性、偶极算符匹配）皆以其为公理基础。
本库原文之一。}
\physics{流生活在半空间 $\mathbb{R}^4\times[0,\infty)$ 上，边界条件
$B(0)=A$；L\"uscher 与 Weisz 将其表述为一个五维局域场论，推导出
Feynman 规则，并用幂计数加 BRST 对称性证明：严格正流时间下流场的
关联函数\emph{不需要底层四维理论之外的任何重整化}。证明把发散组织成
本身即为重整化后的低维算符之流的反项，系数递归固定；BRST Ward 恒等式
消灭危险的类别。这就是每个"小流时间展开"所依赖的全阶定理：
\begin{equation}
  O_R(x) \;\xrightarrow{t\to0}\; c_O(t,\mu)\,[O]_R(x) + \mathcal{O}(t),
\end{equation}
其中 Wilson 系数 $c_O$ 可计算，极限有限而无歧义。}
\impact{本书第三、六章中每一个匹配观测量——跑耦合、能动张量、手征
凝聚、涂抹 twist-2 算符、涂抹准 PDF——都是该展开的具体实例，其系数
都在此文引入的五维形式体系内计算。}
\cardend

% ----------------------------------------------------------------------
\section{费米子与手征对称性}
\papercard{M.\ L\"uscher}{手征对称性与 Yang--Mills 梯度流}
{JHEP \textbf{04} (2013) 123}
{arXiv:1302.5246 | DOI: 10.1007/JHEP04(2013)123}
\whyselected{流语料第三大枢纽（约 $\sim$12/99）：此处定义的费米子流是
ringed fermion、涂抹双线性、涂抹准 PDF 纲领及库内流化 CP 破坏算符
的共同基础。本库原文之一。}
\physics{通过协变热方程把流推广到夸克场
\begin{equation}
  \partial_t\chi = \Delta\,\chi, \qquad \Delta = D_\mu D_\mu,
  \qquad \chi(0)=\psi ,
\end{equation}
其解容许一个五维局域理论的表示（费米子在额外一维中传播）。本文证明
或构造的结果：流双线性如 $S(t)=\bar\chi\chi$ 具有有限连续极限；
$S(t)$ 的小流时间极限定义重整化标量密度，故手征凝聚来自
\begin{equation}
  \lim_{t\to0} Z_S^{-1}(t)\,\langle S(t)\rangle = -\langle\bar\psi\psi\rangle_R ,
\end{equation}
没有加性自重整化的歧义；此外可在格点作用量中加入流诱导的维度五项
构造 O($a$) 改进——这正是后来用于保持涂抹准 PDF 波函数重整化可乘的
"ringed fermion"思想的胚芽。}
\impact{Monahan--Orginos 的涂抹准 PDF（第六章）与 Shindler 的梯度流矩
方案（库内）都从流费米子出发；库内的 NNLO 双线性匹配论文计算的正是
该展开的系数。}
\cardend

% ----------------------------------------------------------------------
\section{从流得到能动张量}
\papercard{H.\ Suzuki}{由 Yang--Mills 梯度流得到能量动量张量}
{Prog.\ Theor.\ Exp.\ Phys.\ \textbf{2013}, 083B03}
{arXiv:1304.0533 | DOI: 10.1093/ptep/ptt059}
\whyselected{流语料的枢纽之一（约 $\sim$8/99），与其姊妹篇 Makino--
Suzuki, PTEP \textbf{2014}, 063B02 [arXiv:1403.4772]（把公式推广到含
费米子的完整 QCD，约 $\sim$9/99）成对出现。两者均为本库原文。}
\physics{既然流的乘积紫外有限（上文式~\ref{eq:flow} 与 L\"uscher--
Weisz 定理），任何恰当归一的流场局域乘积都有无歧义的小流时间展开。
Suzuki 构造了小流时间展开以无迹应力张量为首项的规范不变组合
$U_{\mu\nu}(t,x)$，并对纯 YM 导出主公式：
\begin{equation}
  T_{\mu\nu}(x)
   = \lim_{t\to0}\Big[\mathcal{C}_1(t)\,U_{\mu\nu}^{(1)}(t,x)
      + \mathcal{C}_2(t)\,\delta_{\mu\nu}\,U^{(2)}(t,x)\Big],
\end{equation}
系数通过迹恒等式 $T^\mu{}_\mu=(\beta(g)/2g)\,F^2$ 与维数规整匹配而
完全固定。结果架起格点与连续方案的桥梁，把迹反常、熵密度等热力学量
变成流的测量。}
\impact{Makino--Suzuki 推广提供了库内面向热力学的流论文所需的费米子
项；Harlander--Kluth--Lange（库内）把 Wilson 系数推至 NNLO，使方法
在定量上具备竞争力。在概念上，这篇论文是\emph{任何}流复合观测量之
模板：先在有限 $t$ 定义，再展开、匹配。第六章将同一模板应用于部分子
物理。}
\cardend
